\chapter{Literature Review}
\label{chap:literature-review}

\section{Social Media Marketing and Digital Content Marketing}
\label{sec:2-1-social-media-marketing-and-digital-content-marketing}

Social media platforms have transformed marketing communication from a one-directional broadcasting activity into an interactive and networked process. In traditional marketing, firms communicated with consumers largely through controlled channels such as print advertising, television, outdoor promotion, or direct sales. Social media, by contrast, functions as a hybrid element of the promotion mix: it allows firms to communicate with customers while also enabling customers to communicate directly with one another. Marketing content can be circulated, interpreted, and redistributed by users themselves \parencite{kaplan2010,mangold2009}. This shift changes both the mechanism and the measurement of marketing effectiveness. Firms must consider not only whether a message is delivered, but also whether users read it, spend time with it, share it with others, and allow it to travel through interpersonal networks.

The practical significance of social media marketing is reflected in the increasing role of social platforms in brand communication and consumer decision-making \parencite{curalate2017,lamontagne2015,pwc2018,read2018}. Social platforms provide firms with channels for disseminating promotional content, maintaining customer relationships, and encouraging user participation. They also allow marketing messages to move beyond a firm's direct audience through sharing, commenting, liking, and other forms of user interaction. This feature is especially important in agent-based social marketing, where a firm does not simply publish content to a public page but relies on individual employees or agents to introduce content into their own interpersonal networks. In such a setting, marketing effectiveness emerges from the interaction between firm-generated content, the agent who shares it, and the audience that receives or redistributes it.

Digital content marketing provides an important conceptual foundation for this dissertation. \textcite{hollebeek2019} define digital content marketing as the creation and dissemination of relevant and valuable brand-related content through digital platforms, with the aim of fostering customer engagement, trust, and relationships instead of focusing only on direct purchase persuasion. This definition is useful for the present study because it shifts attention from immediate sales conversion to the broader relational effects of marketing content. Marketing articles shared through WeChat are not merely advertisements; they are content objects that may attract attention, communicate expertise, sustain reader interest, and encourage further interaction with the firm or its representatives. Their effectiveness is partly associated with whether readers perceive the content as relevant, valuable, and credible.

The concept of customer engagement further clarifies why social media marketing should be studied as an interactive process. \textcite{brodie2011} characterise customer engagement as a relational concept rooted in interactive experience and value co-creation, in which customers are not passive recipients of promotional information but active participants who interpret, respond to, and sometimes redistribute brand-related content. This perspective is especially relevant to WeChat-based marketing. When a reader opens an article, spends time reading it, or reshares it to others, the reader participates in the diffusion process. Diffusion effectiveness is not limited to the initial delivery of content. It is reflected in a sequence of user behaviours, including exposure, reading, continued attention, and further sharing.

Trust is also central to social media marketing. Compared with anonymous advertising placements, messages shared through interpersonal networks may be received differently because they are attached to a familiar or identifiable source. Source-credibility research suggests that the persuasiveness of a communication depends partly on the perceived expertise and trustworthiness of the source \parencite{ohanian1990}. In agent-based marketing, the marketing agent is not merely a technical channel for distributing content. The agent is also an organisational representative whose identity, role, and perceived expertise may shape how the message is interpreted. This consideration is particularly relevant in industries such as interior design, home decoration, and real-estate-related services, where trust, professional knowledge, and perceived relevance are likely to shape how consumers evaluate marketing information.

At the same time, the networked character of social media marketing creates a practical challenge. Content is not equally relevant to all agents or all audience groups. A one-size-fits-all strategy, under which the same article is distributed broadly to many agents without considering their professional role, audience context, or topical relevance, may lead to inefficient diffusion. It may also waste social capacity: agents have limited time, credibility, attention, and willingness to share promotional content within their personal networks. If agents repeatedly share content that is weakly related to their role or audience, readers may become less responsive and agents may become less willing to participate in future marketing activities.

This stream of literature provides the conceptual background for this dissertation. Social media marketing and digital content marketing show that marketing effectiveness depends on more than the existence of a communication channel. It also depends on whether content attracts attention, builds trust, encourages engagement, and moves through interpersonal relationships. In the WeiMai WeChat marketing context examined in this study, these issues appear in an agent-based form: promotional articles are shared by marketing agents through their personal networks, and readers' opening, reading, and resharing behaviours form observable diffusion records.

\section{Message Content and Online Diffusion}
\label{sec:2-2-message-content-and-online-diffusion}

The effectiveness of social media marketing rests partly on the content that is communicated and partly on the agent-mediated social context through which it is distributed. The first perspective treats the message itself as a central factor in online word of mouth and social media diffusion, and seeks to explain why some messages are read, discussed, and shared while others are ignored.

\textcite{berger2011} examine the psychological drivers of immediate and ongoing word of mouth, a behaviour closely related to product diffusion and marketing communication. Drawing on daily conversation data across multiple products, together with field and laboratory experiments, they show that content characteristics and contextual cues are associated with whether people talk about products and ideas. Their study is important for the present dissertation because it shifts attention from the existence of a network to the properties of the message being transmitted. A network may provide possible pathways for diffusion, but the content itself helps determine whether audiences have a reason to read, discuss, or share a message.

Research on emotional content further reinforces the importance of message characteristics. \textcite{heath2001} study emotional selection in the spread of urban legends and show that emotionally evocative content is more likely to be transmitted. \textcite{bergermilkman2012} develop this logic in an online news context, finding that virality is associated not only with whether content is positive or negative, but also with the level of emotional arousal it generates. High-arousal emotions, such as awe, anger, and anxiety, are more strongly associated with sharing than low-arousal emotions. These findings suggest that users do not share content only because it is informative. They may also share it because it attracts attention, activates emotion, and supports social expression.

Other studies examine observable content features in social media environments more directly. \textcite{lahuertaotero2018}, for example, analyse 30,082 tweets from 10,120 Twitter users and identify several message-level features associated with diffusion and electronic word of mouth. Their findings specifically identify hashtags, mentions, links, sentiment, and tweet length as tweet elements that influence diffusion and popularity. This evidence is relevant to the present study because WeChat marketing articles also vary in topic, length, visual richness, and presentation style. Although this dissertation does not attempt to measure every linguistic or psychological dimension of article content, it follows the same general premise that observable content characteristics may be systematically associated with diffusion outcomes.

In the context of article-based marketing, topic relevance is particularly important. A marketing article is not only a neutral container of information, it signals what type of problem, product, service, or lifestyle issue the firm is asking readers to attend to. In an interior-design marketing context, for example, articles may focus on home decoration, renovation design, real estate, promotional events, brand information, lifestyle culture, or customer service. These topics may differ in how easily they attract attention, how directly they connect to readers' current needs, and how suitable they are for resharing within personal networks. Topic category provides one way to represent substantive content differences in empirical diffusion analysis.

The coherence between different parts of a message may also matter. In online articles, the title usually provides the first cue that shapes reader expectations, while the body content determines whether those expectations are sustained. When the title and body communicate a consistent topic, readers may experience the article as more coherent and relevant. When the title suggests one topic but the body content develops another, the mismatch may reduce perceived relevance, shorten attention, or discourage resharing. In this dissertation, this idea is operationalised through a title–content semantic consistency measure.

Visual and structural features of marketing articles may also shape user engagement. Images may make an article more concrete, visually salient, and easier to process, especially in product- or design-oriented industries where visual presentation is closely tied to perceived value. Word count may also matter because it reflects the amount of information provided and the effort required from readers. A very short article may provide insufficient detail, while a very long article may require more attention than readers are willing to invest. These features do not determine diffusion on their own, but they provide measurable indicators of article presentation that can be examined alongside topic and semantic consistency.

Taken together, the message-content literature motivates the first level of the empirical analysis in this dissertation. At the content level, the study examines whether article characteristics are associated with diffusion outcomes across article–agent diffusion cases. The main content variables include topic category, title–content semantic consistency, word count, and image-use indicators. These variables do not exhaust all possible dimensions of content quality, but they provide a structured and measurable representation of marketing articles suitable for empirical analysis.

At the same time, content-based explanations remain incomplete. The same article may diffuse differently when shared by different agents, and the same agent may show different outcomes for different types of articles. This means that article content alone cannot fully explain diffusion effectiveness in an agent-based social marketing system.

\section{Network Structure, Agents, and Diffusion}
\label{sec:2-3-network-structure-agents-and-diffusion}

In an agent-based social marketing system, the firm does not simply publish an article to an anonymous audience. The article is shared by a particular marketing agent through that agent's WeChat contacts. The agent is therefore the entry point of the observed diffusion cascade, and agent-level variation is relevant because different agents connect the same firm-generated content to different interpersonal contexts.

\textcite{katz1955}'s two-step flow account argues that information often reaches wider audiences through individuals who interpret and relay messages within their social circles. It establishes a durable communication logic: messages often travel through identifiable intermediaries, and these intermediaries affect how information enters interpersonal audiences.

Online diffusion research restates this logic in network terms. \textcite{katona2011} examine the growth of an invitation-based online social network and model how network exposure affects individual adoption. Their results distinguish between a degree-related mechanism, in which individuals are more likely to adopt when they are connected to more adopters, and a clustering-related mechanism, in which adoption is affected by how densely connected the adopting neighbours are. This distinction is useful because it shows that diffusion is shaped not only by whether a message or product is present in a network, but also by the local structure through which exposure occurs.

\textcite{yoganarasimhan2012} reaches a related conclusion in a different empirical setting. In a study of YouTube video propagation, video popularity is linked not only to video content but also to the local network structure of the author or seed user. This finding is relevant to the present dissertation because marketing agents play a comparable structural role at the beginning of each observed cascade. They are not necessarily influential in the public-platform sense, but they are the points through which articles first enter distinct interpersonal networks. Taken together, these studies do not imply that agent characteristics mechanically determine diffusion. They do show why the sender and the network context around the sender cannot be ignored.

For the present dissertation, this literature motivates an agent-characteristic level of analysis. The relevant observable characteristics are limited: job category, gender, and the number of article-sharing cases observed for each agent. Job category is used as a coarse indicator of professional context, gender is included as an observable demographic attribute, and article count is used as an exposure and stability control. These variables do not directly observe each agent's full audience network or recipient preferences. Their role is more modest: they allow the analysis to test whether observable differences between agents are associated with differences in realised diffusion outcomes.

The purpose of this section is consequently narrow. It establishes why agents are meaningful units of analysis in this setting and why the empirical design separates observable agent characteristics from reconstructed diffusion-network outcomes.

\section{Content–User Fit and Content–Agent Fit}
\label{sec:2-4-content-user-fit-and-content-agent-fit}

Message content and the agent-mediated social context are two important perspectives on online diffusion. However, neither perspective fully explains a recurring pattern in agent-based marketing: the same article may diffuse strongly through some agents but weakly through others, while the same agent may generate different outcomes for different types of articles. A message may be well written and visually attractive, yet still perform poorly if it is shared by an agent whose audience has little interest in the topic. Conversely, an agent who generally achieves strong diffusion will not necessarily be equally suitable for every article. This pattern points to a third perspective: the correspondence between what is communicated and the social context through which it is communicated. In this dissertation, this correspondence is discussed as content–user fit when the focus is the match between content and an audience, and as content–agent fit when the focus is the match between article topics and the sharing agent's professional context.

A theoretical basis for why such fit may matter is provided by the principle of homophily. \textcite{mcpherson2001} show that social ties tend to form disproportionately among people who are similar in characteristics such as occupation, interests, values, and demographic background. In diffusion settings, this implies that an agent's observable attributes may carry information about the audience the agent is likely to reach. An agent whose professional role is closely related to interior design, for example, may have contacts who are more receptive to home decoration or renovation-related content than to unrelated promotional material. Homophily provides a theoretical rationale for treating agent characteristics as imperfect but useful proxies for audience context. It also suggests that the alignment between article topic and agent context may be associated with how readily a message is read, interpreted, and reshared.

At the same time, homophily also cautions against strong causal claims in observational network research. \textcite{shalizi2011} argue that homophily and social contagion are generally difficult to separate in observational social network data, because similar individuals may both connect with one another and behave similarly for reasons unrelated to direct influence. This point is important for the present dissertation. The study does not claim that content–agent fit causally produces diffusion outcomes. Instead, it examines whether stronger fit is systematically associated with stronger observed diffusion performance, while interpreting the results as explanatory and associational, not causal.

Empirical research on referral and interpersonal marketing also supports the relevance of fit. \textcite{vandenbulte2018} examine a customer referral programme and show that customers acquired through referrals can exhibit higher value and lower churn than customers acquired through other channels. One explanation is that referrals may involve better matching between products and potential customers, because existing customers pass information to people for whom the offering is more relevant. Although referral programmes differ from WeChat article sharing, both settings involve interpersonal transmission and social filtering. A message is not sent to an anonymous mass audience. It passes through a person whose social context shapes who receives it and how it is interpreted.

The role of fit is developed more directly in the social media diffusion literature. \textcite{zhang2017} observe that previous research often examined message content and personal influence separately, and they model message content, personal influence, and content–user fit jointly in the context of Twitter rebroadcasting. Their study shows that content–user fit is associated with users' rebroadcasting behaviour, over and above content and influence considered separately. This finding is especially relevant because it suggests that diffusion is related not simply to the message itself or to the general influence of the sender, but also to the match between a message and the audience context.

A related stream of influencer marketing research reaches a similar conclusion through the idea of congruence. \textcite{belanche2021} show that congruence among influencer, product, and consumer is associated with consumer responses, while \textcite{breves2019} find that perceived fit between an influencer and an endorsed brand can strengthen source credibility and persuasive effectiveness. These studies indicate that marketing communication may be more effective when the message, the source, and the audience are perceived as compatible. For the present dissertation, this literature is useful because it shifts attention from sender prominence alone to the relevance of a specific message to a specific sharing context.

This dissertation extends the content–user fit logic to a firm-side WeChat marketing setting. In \textcite{zhang2017}, fit is defined mainly in relation to individual users' historical content preferences and rebroadcasting behaviour. In influencer marketing, fit is often defined as congruence among an influencer, a product or brand, and consumers. In the present study, the focus is different. The marketing agent is not a public influencer and is not simply an ordinary platform user. The agent is an organisational representative, the starting point of a diffusion cascade, and the observable entry point through which firm-generated content enters an interpersonal audience network. The relevant matching relationship is between article topics and the agent context through which the article is distributed. Because homophily implies that an agent's attributes carry information about the agent's audience, fit can be operationalised through the correspondence between article topics and agent characteristics, without requiring direct observation of each recipient's preferences.

Professional role provides one practical way to operationalise this context. In an interior-design marketing system, agents may differ in job function, department, client contact, expertise, and audience expectations. An article about home decoration or renovation design may fit better with agents whose roles are closer to design consultation or customer-facing home services, while an article about promotional events or brand information may fit different types of agents. This does not mean that job category perfectly captures audience preferences. It provides an observable proxy for the professional and social context in which sharing occurs. Content–agent fit is defined here as a measurable approximation of how well an article topic aligns with the agent context through which it is shared.

This framing has direct implications for personalised content assignment. A one-size-fits-all strategy treats agents as interchangeable broadcasting channels and distributes the same content broadly, regardless of topic relevance. The fit perspective suggests a different approach: firms may improve marketing efficiency by assigning articles to agents whose professional context is more closely aligned with the article topic. Such assignment may help conserve social capacity by reducing the circulation of poorly matched content and by increasing the likelihood that shared articles are perceived as relevant by the receiving audience.

Taken together, the literature on homophily, referral matching, social media rebroadcasting, and influencer congruence provides the conceptual basis for the Level 3 analysis in this dissertation. The content model asks whether certain article characteristics are associated with diffusion outcomes. The agent-characteristic model asks whether observable agent attributes relate to agent-level diffusion outcomes. The content–agent matching model asks whether the same type of content is associated with stronger agent–topic diffusion performance when paired with more relevant agent contexts. Consistent with the multi-dimensional view of diffusion effectiveness adopted in this dissertation, this matching relationship is examined against agent–topic reach as the primary outcome and against supplementary diffusion measures where they carry a coherent agent–topic-level interpretation. This third analytical level provides the most direct test of the dissertation's central proposition: in agent-based social marketing, one size does not fit all.

\section{Research Gap and Conceptual Framework}
\label{sec:2-5-research-gap-and-conceptual-framework}

The literature reviewed in this chapter establishes three main findings that frame the present study. First, message content is systematically related to online diffusion. Prior studies show that content characteristics such as interest, emotional appeal, format, and message length are associated with whether messages are read, discussed, and reshared \parencite{berger2011,bergermilkman2012,heath2001,lahuertaotero2018}. Second, the actor who shares a message and the interpersonal context surrounding that actor also matter for diffusion. Research on opinion leadership, network structure, and local network position suggests that information often enters a network through identifiable actors whose audience access shapes how messages spread \parencite{katz1955,katona2011,yoganarasimhan2012}. Third, the fit between content and the receiving or sharing context is associated with communication effectiveness. Studies on referral matching, content–user fit, and influencer congruence indicate that messages may perform better when they are aligned with the people and contexts through which they are transmitted \parencite{belanche2021,breves2019,vandenbulte2018,zhang2017}.

The gap addressed by this dissertation is therefore narrow. The less developed issue is how content–context fit should be defined and tested in a firm-side, agent-mediated WeChat marketing system, where organisational agents distribute firm-generated articles through their own interpersonal networks. The practical decision in such a system is not simply whether an individual user will rebroadcast a message, nor whether a public influencer is congruent with a brand. It is whether a particular article topic is better assigned to one type of marketing agent than to another.

This setting makes content–agent fit more specific than general content–user fit. In \textcite{zhang2017}, fit is measured in relation to users' content preferences and rebroadcasting behaviour. In influencer marketing, congruence is usually considered among an influencer, a product or brand, and consumers \parencite{belanche2021,breves2019}. In referral research, matching is connected to the social filtering that occurs when existing customers introduce an offering to potential customers \parencite{vandenbulte2018}. These studies support the broader relevance of fit, but they do not specify how to measure fit when the observable matching unit is an article topic paired with an organisational agent's professional context. This is the main conceptual and empirical gap that motivates the Level 3 analysis.

A related methodological implication follows from the same setting. Diffusion effectiveness is not exhausted by reach. A cascade may reach many readers, remain shallow, spread through several layers, prompt resharing, or become concentrated around particular nodes. Measures such as structural virality were developed to describe aspects of diffusion shape that simple popularity counts cannot capture \parencite{goel2016}. In the present study, however, these cascade and network measures are reconstructed after diffusion has occurred. They are therefore useful as supplementary outcomes or descriptors, but they should not be treated as pre-diffusion predictors of the same process from which they are calculated.

In response, this dissertation adopts a three-level conceptual framework. At the content level, the study examines whether article characteristics are associated with diffusion outcomes across article–agent cases. This level draws on the message-content literature and focuses on topic category, title–content semantic consistency, word count, and image-use indicators. At the agent-characteristic level, the study examines whether observable agent characteristics are associated with agent-level diffusion patterns. This level draws on the agent and network literature and focuses on job category and gender, with agent activity included as a control. Reconstructed diffusion-network measures are treated as agent-level outcomes or descriptors. At the content–agent matching level, the study examines whether alignment between article topics and agent professional contexts is associated with stronger agent–topic diffusion performance.

Two principles guide this framework. The first is the separation of units of analysis. Content characteristics, agent characteristics, and content–agent fit are observed at different empirical levels: article–agent diffusion cases, agents, and agent–topic combinations. Modelling them separately keeps coefficient interpretation aligned with the level at which each variable is defined. The second principle is leakage-aware outcome treatment. Reach is used as the primary indicator of realised diffusion performance, while reconstructed cascade and network metrics are used only where they have a coherent outcome or descriptor interpretation. Homophily provides the theoretical rationale for treating observable agent characteristics as imperfect proxies for audience context. At the same time, the observational nature of the data requires the findings to be interpreted as explanatory associations, not causal effects.
